% ==============================================================================
% Table 2: Biochemical Parameters of Botanical Defenses and Signals
% Placed at bottom [b] of page, caption above table
% ==============================================================================
%
% NARRATIVE STORY OF TABLE 2:
% This table establishes the empirical calibration underpinning the biological
% simulation engine in PHIDS. It bridges established empirical literature in plant
% physiology and chemical ecology with the formal computational parameters used across
% the coupled hybrid system.
%
% The narrative progresses through five functional tiers of plant defence:
% 1. Constitutive Morphological Barriers: Static structures (trichomes, biogenic
%    silica, lignified secondary cell walls) that require upfront energetic allocation
%    during tissue maturation but impose zero ongoing synthesis delay upon attack.
% 2. Airborne VOC Signals (Green Leaf Volatiles & Terpenoids): Rapidly emitted semiochemical
%    plumes that diffuse through the turbulent atmospheric boundary layer to prime downwind
%    conspecifics, characterised by high dispersion and finite photochemical half-lives.
% 3. Subterranean Mycorrhizal Networks (CMNs): Underground hyphal networks that provide
%    shielded, out-of-band communication conduits connecting root systems, operating at a
%    continuous metabolic maintenance fee (carbon tax) paid by the host plant.
% 4. Inducible Allochemical Defenses (Glucosinolates, Alkaloids, Proteinase Inhibitors):
%    Metabolically expensive secondary toxins synthesized de novo upon tissue wounding,
%    subject to a critical physiological lag phase (synthesis delay theta_synth).
% 5. Vascular Phloem Translocation (Carbon/Nitrogen Sinks): Rapid enzymatic reallocation
%    of mobile sugars and amino acids away from grazed foliar tissue into subterranean root
%    storage, dynamically reducing apparent tissue nutrition factor eta(t).
%
% ==============================================================================

\begin{table*}[b]
\caption{Biochemical, Physiological, and Kinetic Parameters Governing Botanical Defence Cascades in PHIDS.}
\label{tab:biochemical_parameters}
\vspace{4pt}
\centering
\footnotesize
\begin{tabularx}{\textwidth}{
  >{\raggedright\hsize=0.90\hsize}X
  >{\raggedright\hsize=1.00\hsize}X
  >{\centering\hsize=0.95\hsize}X
  >{\centering\hsize=0.80\hsize}X
  >{\centering\hsize=0.90\hsize}X
  >{\raggedright\hsize=1.15\hsize}X
  >{\raggedright\arraybackslash\hsize=1.30\hsize}X
}
\toprule
\textbf{Defence Modality} &
\textbf{Biochemical Basis} &
\textbf{Metabolic Cost} &
\textbf{Latency ($\bm{\theta}$)} &
\textbf{Decay / Half-Life} &
\textbf{Empirical Basis} &
\textbf{Model Symbol / Role} \\
\midrule

\textbf{Constitutive Armour} &
Lignin, cellulose, silica trichomes &
Continuous fraction \newline $\delta_{\mathrm{morph}} \in [0, 0.50]$ &
$0\,\text{ticks}$ (Permanent) &
$\tau_{1/2} \to \infty$ (Permanent) &
Gatehouse \cite{gatehouse2002plant} &
$\delta_{\mathrm{morph}}$ (digestibility), \newline
$c_{\mathrm{bite}}$ (mechanical wear) \\

\textbf{Airborne VOC Plumes} &
$(E)$-$\beta$-ocimene, green leaf volatiles &
Signal emission $Q_{\mathrm{emit}}$; zero resting tax &
$1\text{--}3\,\text{ticks}$ (Plume onset) &
Clearance $\gamma_s \approx 1 - \lambda_s\Delta t$ (diffusion $\sigma = 0.40$) &
Farmer \cite{farmer2001surface} &
$Q_{\mathrm{emit}}$ (emission flux), \newline
$\gamma_s$ (clearance retention) \\

\textbf{Mycelial Conduction} &
Hyphal action potentials, jasmonates &
Link upkeep $\delta_E^{\mathrm{tax}}$ / link &
$1\,\text{tick/hop}$ ($1\text{--}3\,\text{ticks}$) &
Unattenuated across cardinal links &
Babikova et al. \cite{babikova2013underground} &
$\delta_E^{\mathrm{tax}}$ (link carbon tax), \newline
$g_{\mathrm{relay}}$ (relay gain) \\

\textbf{Inducible Toxins} &
Glucosinolates, nicotine, proteinase inhibitors &
Active synthesis tax $c_{\mathrm{synth}}$ / tick &
$8\text{--}12\,\text{ticks}$ ($1\text{--}2$ in compact orchard runs; $\theta_{\mathrm{synth}}$) &
Active persistence $\tau_{\mathrm{aft}}$ ticks &
Gatehouse \cite{gatehouse2002plant} &
$\theta_{\mathrm{synth}}$ (induction lag), \newline
$\mu_{\mathrm{tox}}$ (lethality rate) \\

\textbf{Primed Acceleration} &
Pre-transcribed mRNA, inactive enzyme pools &
$\approx 1\%$ resting alert upkeep &
$1\text{--}3\,\text{ticks}$ ($\theta_{\mathrm{primed}}$) &
Gated by Hill saturation $S(C) \ge \tau_{\mathrm{prim}}$ &
Gatehouse \cite{gatehouse2002plant} &
$n$ (Hill exponent), \newline
$K_d$ (half-saturation) \\

\textbf{Phloem Translocation} &
Apparent foliar nutrition withdrawal &
Stress withdrawal without biomass loss &
Rate-limited recovery ($\gamma_{\text{trans}}$) &
Relaxation: $\eta(t+1) = \eta(t) + \Delta \eta$ &
Schwachtje et al. \cite{schwachtje2006snf1} &
$\gamma_{\mathrm{trans}}$ (recovery rate), \newline
$\eta(t)$ (palatability factor) \\

\bottomrule
\end{tabularx}
\end{table*}
